\subsection{Requirements and Intent Extraction, Traceability, and Natural-Language-to-Formal Translation}
\label{subsec:rw9}

A pipeline that treats a machine-checkable contract as the durable source of truth must first answer where that contract comes from. In brownfield settings the contract is not authored from a blank page but \emph{recovered} from three competing streams of evidence---code, issue tickets, and documentation---each of which a substantial literature has studied in isolation. This subsection surveys that literature thematically and positions the gap that the present article addresses.

\paragraph{Traceability link recovery.}
The foundational problem of relating informal artefacts to code was framed by Antoniol et al., who recovered links between source code and free-text documentation using information-retrieval (IR) models on the premise that programmers embed meaningful names in identifiers \cite{antoniol2002recovering}. De Lucia et al. extended this line into an integrated artefact-management environment using latent semantic indexing, demonstrating that IR-based recovery could be embedded in a working development tool rather than run as a one-off batch analysis \cite{delucia2007recovering}. A decade of such work was consolidated by Borg et al., whose systematic mapping catalogued the dominant vector-space and topic-model techniques and, tellingly, the persistent precision--recall ceiling that purely textual similarity imposes \cite{borg2014recovering}. The field's research agenda was crystallized by Gotel et al.'s grand-challenge roadmap, which named ubiquitous, trustworthy, and \emph{purposed} traceability as the long-horizon goal and observed that links created without a downstream consumer are rarely maintained \cite{gotel2012grand}. More recent work has displaced shallow IR with learned semantics: Guo et al. applied recurrent architectures with domain word embeddings to inject artefact semantics into the tracing decision \cite{guo2017semantic}, and Lin et al.\ showed that pretrained BERT models (T-BERT) transfer effectively to the low-resource trace-link setting, materially improving link accuracy over IR baselines \cite{lin2021tbert}. Across this body of work, however, a recovered link asserts only \emph{relatedness}; it does not commit to any machine-checkable claim about behaviour, and so cannot by itself serve as a verification oracle.

\paragraph{Mining requirements and intent from issue trackers.}
Issue trackers are now the de facto repository of just-in-time requirements. Merten et al.'s grounded analysis of open-source trackers established that most requirements-relevant content---feature clarification, edge cases, rationale, and bug-driven constraints---is buried in unstructured discussion rather than in structured fields \cite{merten2015issuetracking}. The scale and noisiness of this evidence stream are well documented by curated corpora such as the public Jira dataset of Montgomery et al., spanning millions of issues, comments, and inconsistently labelled links \cite{montgomery2022jira}. This work substantiates two assumptions of the present pipeline's characterisation phase: that tickets are a genuine source of \emph{why}---the intent and known edge cases that code alone cannot reveal---and that linking tickets to code is inherently heuristic and noisy, credible only on a curated module rather than across a whole legacy application.

\paragraph{Documentation--code consistency.}
A parallel thread treats documentation not as a source to be mined but as a claim to be checked against code. Tan et al.'s \texttt{@tComment} inferred null/exception properties from Javadoc and generated tests to expose comment--code inconsistencies, surfacing defects in either the comment or the body \cite{tan2012tcomment}. Wen et al.'s large-scale study confirmed that such inconsistencies are common, persistent, and frequently co-located with bugs \cite{wen2019codecomment}. This is precisely the disagreement-detection value the present article elevates from a one-off analysis to a first-class lifecycle signal: a divergence between code and its documentation is a \emph{candidate bug}, not an error to be silently reconciled.

\paragraph{Natural-language-to-formal translation.}
Bridging informal intent to a checkable artefact is the closest prior art to the elaboration layer. Cosler et al.'s \texttt{nl2spec} translated unstructured English to temporal logic with LLMs and---crucially---introduced an interactive loop in which sub-formulas are mapped back to natural-language fragments for human correction \cite{cosler2023nl2spec}. Endres et al. studied whether LLMs can turn natural-language intent into formal method postconditions, validating the generated contracts by their power to discriminate buggy from correct code and catching dozens of real Defects4J faults \cite{endres2024nl2postcond}. These results directly motivate the present pipeline---inference can remove the spec-authoring bootstrap cost---but they also expose the residual risk it must manage: LLM-produced formal text is sometimes wrong or vacuous, so it cannot be a trusted oracle absent independent authority and human ratification.

\paragraph{The gap.}
Each thread above optimises one edge of the intent problem---link recovery relates artefacts, issue mining extracts intent, consistency checking flags divergence, and NL-to-formal translation produces a candidate spec---but none composes them into a single, provenance-tagged, machine-checkable contract that serves simultaneously as the agent's input and its verification oracle. None reconciles the three evidence streams into a contract whose clauses carry per-source provenance and graded assurance, nor treats inter-stream disagreement as the deliberate success outcome of a brownfield startup phase. Critically, prior NL-to-formal work validates a spec against a fixed code sample, whereas the present article verifies code against the recovered contract over \emph{all} inputs via OpenJML extended static checking, and propagates a callee-spec change to recomputed caller obligations (building on Article~3 of this programme \cite{edmonds_article3}). The contribution surveyed against here is therefore not a better link recovery or translation technique, but the architectural integration that turns recovered intent into a verifiable, compositional contract layer.
